\chapter{René Frédéric Thom}\label{ch1958renethom:renuxe9-fruxe9duxe9ric-thom-fields-medal-1958-a-comprehensive-report}

\begin{center}\rule{0.5\linewidth}{0.5pt}\end{center}

\textbf{Prepared:} July 2026
\textbf{Subject:} René Frédéric Thom (2 September 1923 -- 25 October 2002)
\textbf{Position:} Professor, Institut des Hautes Études Scientifiques (IHÉS)
\textbf{Doctoral Advisor:} Henri Cartan (Université de Paris / École Normale Supérieure)
\textbf{Fields Medal Citation:} \emph{``The creator of cobordism theory \cite{thom1954}, a central
concept in algebraic topology that classifies manifolds up to the existence of a
common boundary.''}

\begin{center}\rule{0.5\linewidth}{0.5pt}\end{center}

\begin{prereqs} Differential topology (manifolds, transversality),
basic singularity theory, multivariable calculus. \end{prereqs}

\begin{objectives} State the classification theorem for smooth manifolds
of dimension $< 5$; explain the cobordism theory; describe the seven
elementary catastrophes. \end{objectives}

\section{Problem History}

The central problem motivating Thom's work was the classification of smooth
manifolds. By the early 1950s, topologists understood homotopy equivalence and
homeomorphism reasonably well, but the finer question of when two manifolds are
\emph{diffeomorphic}---admitting a smooth bijection with smooth
inverse---remained largely unexplored.

A natural approach was to ask which manifolds are boundaries of higher-dimensional
manifolds. Two closed $n$-manifolds $M_1$ and $M_2$ are \textbf{cobordant} if
their disjoint union $M_1 \sqcup M_2$ is the boundary of a compact
$(n+1)$-manifold $W$. The equivalence classes form an abelian group $\Omega_n$.

Pontryagin showed cobordant manifolds must have equal Stiefel--Whitney numbers.
Whether matching numbers guarantees cobordism was the central open problem.

The study of singularities of smooth maps was in its infancy. Whitney showed
generic maps $\mathbb{R}^2 \to \mathbb{R}^2$ have only fold and cusp
singularities, but a general classification was lacking.

The difficulty was both conceptual and technical, requiring new invariants
bridging topology and differential geometry.

Hopf, in his 1958 ICM laudatio, emphasized that Thom unified manifold
classification and singularity analysis through characteristic classes and
transversality.

The algebraic topology community had developed spectral sequences and Steenrod
operations, but these had not been applied systematically to manifold geometry.

Before Thom, classification was known only in low dimensions. The classical genus
classification handled surfaces, but 3-manifolds were already complicated.

Characteristic classes---Stiefel--Whitney, Pontryagin, and Chern---provided raw
material, but their relationship to manifold geometry was not fully understood.

Thom's insight was to use these classes as coordinates on the cobordism ring.

The problem had connections to physics: gauge theory involves infinite-dimensional
moduli spaces whose topology Thom's tools would later illuminate.

\textbf{Historical context}: The cobordism concept originated with Lefschetz
(1926), who showed that a closed oriented surface bounds a three-manifold if and
only if its genus is even. Pontryagin (1950) computed the oriented cobordism
groups $\Omega_n^{\text{SO}}$ for $n \leq 5$, finding $\Omega_0^{\text{SO}} = \mathbb{Z}$,
$\Omega_1^{\text{SO}} = \Omega_2^{\text{SO}} = 0$, $\Omega_3^{\text{SO}} = 0$,
$\Omega_4^{\text{SO}} = \mathbb{Z}$.

The \textbf{unoriented cobordism} problem was more subtle. While Stiefel--Whitney
numbers provided necessary conditions, their sufficiency was unknown. The key
obstruction came from understanding how characteristic classes interact with
boundaries.

\textbf{Whitney's work on singularities} (1955) classified generic maps from
surfaces to $\mathbb{R}^3$, showing that fold and cusp singularities are the
only stable types. However, a general theory for maps between arbitrary manifolds
remained elusive.

\section{How Was It Solved}

Thom's key insight was the \textbf{Pontryagin--Thom construction}. Given
$M^n \subset \mathbb{R}^{n+k}$ with normal bundle classified by
$\nu: M \to BO(k)$, collapsing the complement of a tubular neighborhood yields:

\[S^{n+k} \to T(\gamma^k)\]

This gives $\Omega_n \cong \pi_{n+k}(MO)$, where $MO$ is the Thom spectrum.

\subsection{The Pontryagin--Thom Construction}

Let $M^n$ be a closed smooth manifold embedded in $\mathbb{R}^{n+k}$. The
\textbf{normal bundle} $\nu$ has rank $k$ and is classified by a map
$\nu: M \to BO(k)$. Taking a tubular neighborhood $T(\nu) \subset \mathbb{R}^{n+k}$
and collapsing its complement to a point yields the \textbf{Thom space}:

\[S^{n+k} \cong \mathbb{R}^{n+k} \cup \{\infty\} \to T(\nu) = D(\nu)/S(\nu)\]

The classifying map for $\nu$ composed with the projection gives:

\[\phi: S^{n+k} \to T(\gamma^k) \to MO(k)\]

where $\gamma^k$ is the universal $k$-plane bundle over $BO(k)$. Taking the
colimit over $k$ gives $\phi: S^{n+k} \to MO$, defining an element of
$\pi_{n+k}(MO)$.

\textbf{The Pontryagin--Thom theorem} states that this correspondence induces an
isomorphism $\Omega_n^{\text{O}} \cong \pi_{n+k}(MO(k))$ for sufficiently large
$k$, and passing to the limit gives $\Omega_*^{\text{O}} \cong \pi_*(MO)$.

\subsection{The Thom Space and Thom Isomorphism}

For an oriented rank-$n$ vector bundle $\xi: E \to B$ with zero section
$i: B \to E$, the \textbf{Thom space} is $T(\xi) = D(\xi)/S(\xi)$, the quotient
of the disk bundle by the sphere bundle. The \textbf{Thom isomorphism} is the
fundamental tool:

\[H^q(B) \cong H^{q+n}(D(\xi), S(\xi); \mathbb{Z})\]

given by cup product with the \textbf{Thom class} $u \in H^n(D(\xi), S(\xi))$.
This class restricts to the generator of $H^n(S^n) \cong \mathbb{Z}$ on each
fiber.

\textbf{The Thom class} is characterized by:
\begin{enumerate}
\item $u|_{\text{fiber}}$ generates $H^n(\mathbb{R}^n, \mathbb{R}^n \setminus \{0\})$
\item $i^*(u) = 0$ where $i: B \hookrightarrow E$ is the zero section
\end{enumerate}

Pullback by the zero section gives the \textbf{Euler class} $e(\xi) = i^*(u) \in
H^n(B)$, with $e(M)[M] = \chi(M)$. For the tangent bundle $TM$ of a closed
oriented $n$-manifold, $e(TM) = \chi(M) \cdot [M]^*$ in $H^n(M)$.

\subsection{Cobordism Groups and Characteristic Numbers}

The \textbf{Stiefel--Whitney classes} $w_i(\xi) \in H^i(B; \mathbb{F}_2)$ are
characteristic classes for real vector bundles. For the tangent bundle $TM$ of a
closed $n$-manifold $M$, the \textbf{Stiefel--Whitney numbers} are the
evaluations:

\[\langle w_{i_1}(TM) \smile \cdots \smile w_{i_k}(TM), [M] \rangle \in \mathbb{F}_2\]

where $i_1 + \cdots + i_k = n$.

\textbf{Thom's cobordism theorem} states:
\begin{theorem}
A closed $n$-manifold $M$ is null-cobordant if and only if all its
Stiefel--Whitney numbers vanish.
\end{theorem}

The proof uses the Thom isomorphism and Steenrod operations to compute
$H_*(MO; \mathbb{F}_2)$. Thom showed:

\[H_*(MO; \mathbb{F}_2) \cong MO_* \otimes \mathbb{F}_2 \cong \mathbb{F}_2[x_2, x_4,
x_5, x_6, x_8, \ldots]\]

where $\deg(x_i) = i$ and one generator appears in each degree not of the form
$2^k - 1$.

\textbf{The oriented cobordism ring} was computed by Milnor (1957):

\[\Omega_*^{\text{SO}} \otimes \mathbb{Q} \cong \mathbb{Q}[x_4, x_8, x_{12},
\ldots]\]

generated by complex projective spaces $\mathbb{CP}^{2k}$. The integral structure
is more complicated, with 2-torsion appearing in dimensions divisible by 4.

\textbf{Thom's formula} connects Steenrod operations to Stiefel--Whitney classes:

\[w_i(\xi) = \Phi^{-1}(Sq^i(u))\]

where $\Phi: H^*(B) \to H^{*+n}(D(\xi), S(\xi))$ is the Thom isomorphism and
$Sq^i$ are the Steenrod squares. This is the \textbf{Wu formula}.

\textbf{The Wu classes} $v_k(M) \in H^k(M; \mathbb{F}_2)$ of a closed manifold $M$
are defined by the property that for all $x \in H^{n-k}(M; \mathbb{F}_2)$:
\[Sq^k(x) = v_k(M) \smile x\]
Thom proved that the total Wu class $v = \sum v_k$ satisfies $v = Sq(w)$,
relating Wu classes directly to Stiefel--Whitney classes.

\textbf{Key cobordism groups} computed by Thom include:
\begin{align*}
\Omega_0^{\text{O}} &= \mathbb{F}_2, \quad \Omega_1^{\text{O}} = \mathbb{F}_2, \quad \Omega_2^{\text{O}} = \mathbb{F}_2, \\
\Omega_3^{\text{O}} &= 0, \quad \Omega_4^{\text{O}} = \mathbb{F}_2^3, \quad \Omega_5^{\text{O}} = \mathbb{F}_2^2
\end{align*}
The generator of $\Omega_1^{\text{O}}$ is the circle $S^1$; of $\Omega_2^{\text{O}}$
is the real projective plane $\mathbb{RP}^2$; and $\Omega_3^{\text{O}} = 0$ reflects
that every closed 3-manifold bounds a 4-manifold.

\textbf{The cobordism ring structure} was determined by computing the homology of
the Thom spectrum $MO$. Using the Serre spectral sequence for the fibration
$MO(k) \to MO(k+1)$, Thom showed that $H_*(MO; \mathbb{F}_2)$ is a polynomial
algebra. The proof relies on the fact that Steenrod operations preserve the
image of the Thom class under the isomorphism.

\textbf{Oriented cobordism} requires more care. Pontryagin showed that oriented
cobordism classes are determined by Pontryagin numbers and the signature.
Thom proved that $\Omega_*^{\text{SO}} \otimes \mathbb{Q}$ is a polynomial ring
generated by $\mathbb{CP}^{2k}$, while the integral structure involves
$\mathbb{Z}/2$-torsion in dimensions divisible by 4.

\subsection{The Transversality Theorem}

The \textbf{Thom transversality theorem} guarantees that any smooth map can be
approximated by a transverse map. Given a smooth map $f: M \to N$ and a
submanifold $Z \subset N$, $f$ is \textbf{transverse} to $Z$ if:

\[df(T_m M) + T_{f(m)} Z = T_{f(m)} N \quad \forall m \in f^{-1}(Z)\]

\textbf{Thom's transversality theorem}: The set of maps transverse to $Z$ is
dense and open in the $C^\infty$ topology. Moreover, any map is homotopic to a
transverse map.

This result is fundamental to differential topology. It implies that the
preimage $f^{-1}(Z)$ is a smooth submanifold of codimension $\operatorname{codim}(Z)$
when $f \pitchfork Z$. The \textbf{Thom--Porteous formula} expresses the
cohomology class of degeneracy loci in terms of Chern classes.

\textbf{The Thom--Porteous formula} applies to a bundle map $\phi: E \to F$
between vector bundles of ranks $e$ and $f$ over a manifold $M$. The degeneracy
locus $D_k(\phi) = \{m \in M : \operatorname{rank}(\phi_m) \leq k\}$ has
codimension $(e-k)(f-k)$ when nonempty, and its cohomology class is given by
a determinant involving Chern classes of $E$ and $F$.

\subsection{The Thom--Mather Isotopy Theorem}

The \textbf{Thom--Mather isotopy theorem} provides the foundational result
for the theory of singular spaces. Given a proper submersion $f: X \to Y$
between Whitney stratified spaces satisfying the $A_f$ condition, the theorem
asserts that $f$ is a topological fiber bundle. The $A_f$ condition requires
that the restriction of $f$ to each stratum is a submersion and that the
strata satisfy a regularity condition along the fibers.

\textbf{Whitney's conditions} $(A)$ and $(B)$ for stratified spaces were
introduced by Whitney (1965) to ensure that singular spaces have well-behaved
tangent planes. Condition $(A)$ requires that tangent planes to nearby strata
limit to a plane containing the tangent plane of the stratum. Condition $(B)$
is stronger and implies $(A)$.

\textbf{Mather's stability theorem} (1970) classified germs of smooth maps
$\mathbb{R}^n \to \mathbb{R}^p$ up to $\mathcal{A}$-equivalence (right-left
equivalence). For stable maps, the singular set has a stratification by
Boardman symbols $\Sigma^{i_1,i_2,\ldots}$, where each symbol indicates the
corank of successive derivatives.

\subsection{Singularity Classification}

Thom classified the \textbf{simple singularities} $A_k, D_k, E_6, E_7, E_8$---
corresponding to Dynkin diagrams. These are classified by their \textbf{Milnor
number} $\mu$, which counts the number of vanishing cycles.

\textbf{The $A_k$ singularities} are given by $f(x) = x^{k+1}$. The Milnor fiber
has the homotopy type of a wedge of $k$ spheres.

\textbf{The $D_k$ singularities} are $f(x,y) = x^2y + y^{k-1}$. These correspond
to the Dynkin diagram of type $D_k$.

\textbf{The $E_6, E_7, E_8$ singularities} are the exceptional simple
singularities:
\begin{align*}
E_6 &: x^3 + y^4 \\
E_7 &: x^3 + xy^3 \\
E_8 &: x^3 + y^5
\end{align*}

In \emph{Structural Stability and Morphogenesis} (1972), Thom classified the
\textbf{seven elementary catastrophes}: fold, cusp, swallowtail, butterfly, and
three umbilics (elliptic, hyperbolic, and parabolic).

\textbf{The fold catastrophe} is modeled by $V(x) = x^3/3 + ax$, giving a
cusp-shaped equilibrium surface in parameter space.

\textbf{The cusp catastrophe} is $V(x) = x^4/4 + ax^2/2 + bx$, with a cusp
singularity in the bifurcation set.

\textbf{The swallowtail catastrophe} is $V(x) = x^5/5 + ax^3/3 + bx^2/2 + cx$.

\textbf{The butterfly catastrophe} is $V(x) = x^6/6 + ax^4/4 + bx^3/3 + cx^2/2 +
dx$.

\textbf{The elliptic umbilic} is $V(x,y) = x^3 - 3xy^2 + ax^2 + by^2 + cxy$.

\textbf{The hyperbolic umbilic} is $V(x,y) = x^3 + y^3 + axy + bx + cy$.

\textbf{The parabolic umbilic} is $V(x,y) = x^2y + y^3/3 + ax^2 + bxy + cy^2$.

\textbf{The Thom--Mather isotopy theorem} states that proper submersions between
Whitney stratified sets satisfying the Thom $A_f$ condition are topological fiber
bundles.

\section{Impact of This Discovery}

Thom's cobordism theory transformed manifold classification into systematic
algebraic theory. Milnor, Wall, and Novikov completed the picture for all
cobordism theories.

The Pontryagin--Thom construction opened stable homotopy theory as a tool for
manifold classification.

The transversality theorem became the most-used tool in differential topology,
underlying intersection theory, Morse theory, and moduli space construction.

Thom's singularity theory launched an entire field. Arnold extended it to several
variables; Mather proved stability theorems; Boardman classified via symbols
$\Sigma^{i,j,k,\ldots}$.

Catastrophe theory found applications in optics, biology, and engineering, though
some social science applications were controversial.

The Thom--Mather theory of stratified spaces became essential in algebraic
geometry for singular varieties.

The Dold--Thom theorem and Thom--Porteous formula influenced algebraic K-theory
and enumerative geometry.

\textbf{Milnor's work} on exotic spheres (1956) used cobordism methods to show
that $S^7$ admits multiple smooth structures. The group $\Theta_n$ of exotic
spheres is isomorphic to the cobordism group $\Omega_n^{\text{O}}$ for $n \geq 5$.

\textbf{The Atiyah--Hirzebruch spectral sequence} computes cobordism groups from
ordinary cohomology. For any generalized homology theory $h_*$, there is a
spectral sequence with $E^2_{p,q} = H_p(X; h_q(pt))$ converging to
$h_{p+q}(X)$.

\textbf{Conner--Floyd theory} developed complex cobordism $MU_*$ as a universal
cohomology theory for complex manifolds. The Lazard ring $L_* \cong MU_*$
classifies formal group laws.

\textbf{Wu's formula} for Steenrod operations on cobordism classes became
essential for computations. The Wu classes $v_k$ satisfy $Sq^k(x) = v_k \smile
x$ for all $x \in H^{n-k}(M)$.

\section{Modern Views and Improvements}

Cobordism extends to equivariant settings with $\Omega_G^*$ for compact Lie groups,
important in mathematical physics.

The Thom--Mather theorem has been refined into tame topological structures.
Hill--Hopkins--Ravenel used equivariant cobordism for the topological
j-homomorphism.

Modern singularity theory connects to tropical geometry, mirror symmetry, and
Calabi--Yau manifolds via the ADE classification.

The Thom--Boardman stratification remains standard. Active research includes
stratified Morse theory and microlocal sheaf theory.

The Thom spectrum $MO$ continues to be computed in new contexts, connecting to
string theory and topological field theory.

\textbf{Brown--Peterson theory} $BP_*$ is a direct summand of $MU_*$ obtained
by inverting all primes except $p$. It provides a simpler theory for $p$-local
computations.

\textbf{Morava $K$-theory} $K(n)_*$ is a field theory at height $n$. The
chromatic spectral sequence organizes stable homotopy theory in terms of these
heights.

\textbf{The cobordism conjecture} in string theory predicts that every consistent
quantum gravity theory must admit a complete set of boundary conditions, related
to the bordism classification of spacetime defects.

\textbf{Topological quantum field theory} (TQFT) uses cobordism categories as
its foundational structure. The Atiyah--Segal axioms define a TQFT as a
functor from the category of cobordisms to vector spaces.

\textbf{The Landau--Ginzburg model} connects singularity theory to mirror
symmetry. The Milnor ring of a singularity $f$ is isomorphic to the quantum
cohomology of the mirror manifold.

\textbf{Arnold's classification} of singularities using $\mathcal{A}$-equivalence
showed that generic maps between surfaces have only finitely many stability
types.

\textbf{The Thom--Mather theory} of stratified spaces has been extended to
\emph{tame topology} by Gabrielov, Gindikin, and Volkovyskii, providing
rigorous foundations for catastrophe theory applications.

\begin{center}\rule{0.5\linewidth}{0.5pt}\end{center}

\section*{Summary}
\addcontentsline{toc}{section}{Summary}

Thom created cobordism theory, showing that manifolds can be classified by
algebraic invariants computed from characteristic classes, and developed
catastrophe theory as a mathematical framework for studying abrupt changes in
continuous systems. His Pontryagin--Thom construction established the fundamental
bridge between differential topology and stable homotopy theory, while his
work on singularities provided a complete classification of elementary
catastrophes with far-reaching applications in science and engineering.